\chapter{Positionnement de l'alternance et problématique}

\section{Étude du secteur : le compartimentage et le désenfumage}

\subsection{Un secteur méconnu mais vital pour la sécurité}

Le compartimentage et le désenfumage font partie de ces secteurs industriels dont le grand public ignore l'existence, alors qu'ils conditionnent directement la sécurité des bâtiments. Lorsqu'un incendie se déclare, c'est la capacité de l'ouvrage à contenir le feu, à préserver les cheminements d'évacuation et à extraire les fumées qui fait la différence. Une trappe coupe-feu ou un exutoire de désenfumage ne sont pas des accessoires : ce sont des éléments de la chaîne de sécurité, au même titre que les issues de secours.

Cette réalité confère aux solutions de compartimentage et de désenfumage un statut particulier. Elles doivent être conçues avec précision, validées selon des référentiels stricts et déployées dans des contextes projet souvent très variés. Contrairement à des produits totalement standardisés, elles combinent fréquemment une forte exigence réglementaire avec un besoin d'adaptation à des dimensions, à des supports et à des environnements architecturaux spécifiques. La logique de gamme y cohabite donc avec une logique de personnalisation maîtrisée.

Le secteur est également caractérisé par une forte responsabilité technique. Une erreur de configuration, une mauvaise interprétation d'un domaine de validité ou une incohérence entre plans, données et fabrication peuvent avoir des conséquences importantes, non seulement en phase de production, mais aussi au moment de la prescription, de la pose ou de la réception du bâtiment.

\subsection{Un environnement réglementaire exigeant}

Le cadre réglementaire du secteur est particulièrement structurant. Les produits liés à la sécurité incendie sont encadrés par des normes, des essais, des procès-verbaux et des réglementations nationales ou européennes qui définissent des performances à atteindre, des conditions de mise en \oe uvre et des limites d'emploi. Dans ce contexte, la donnée technique n'est pas seulement descriptive : elle est aussi normative. Une valeur de dimension, un type de support, un niveau de classement ou un accessoire autorisé renvoient tous à des conditions de validité qui doivent être respectées.

Cette exigence a des conséquences directes sur le travail d'ingénierie. D'une part, elle impose une traçabilité solide : il faut pouvoir relier une configuration à une source technique, à une règle métier et, si nécessaire, à un document de certification. D'autre part, elle demande une certaine pédagogie interne. Les outils doivent aider les utilisateurs à rester dans un cadre valide sans leur demander de relire systématiquement des documents longs et complexes pour chaque configuration.

Dans les faits, une grande partie de la difficulté ne vient donc pas uniquement de la technique produit, mais de la traduction fidèle d'un corpus réglementaire et documentaire en règles utilisables dans les outils de l'entreprise. C'est précisément là que les enjeux de structuration de données prennent toute leur importance.

\subsection{Les acteurs et orientations du secteur}

Le secteur réunit des industriels spécialisés, des bureaux d'études, des laboratoires d'essais, des prescripteurs, des installateurs et des maîtres d'ouvrage. Chacun intervient avec sa propre logique : l'industriel doit concevoir des solutions fiables et industrialisables ; le bureau d'études cherche la conformité et l'intégration projet ; le laboratoire qualifie les performances ; l'installateur attend des produits clairs à mettre en \oe uvre ; le client final attend une solution sûre, documentée et livrable dans les délais.

Plusieurs orientations structurent aujourd'hui l'évolution du secteur. La première est la transformation numérique des processus. Les outils de configuration, les modèles paramétriques, la gestion documentaire et la traçabilité numérique deviennent des leviers de compétitivité. La deuxième est la montée en puissance des logiques BIM et de continuité numérique, qui poussent les industriels à mieux structurer leurs données et à les rendre réutilisables en dehors du seul cadre interne. La troisième concerne les enjeux environnementaux et RSE : réduction des reprises, limitation du gaspillage, amélioration de la durabilité des produits, meilleure qualité documentaire et meilleure maîtrise des modifications au fil du cycle de vie.

Dans ce paysage, les entreprises capables de transformer leur expertise métier en outils fiables disposent d'un avantage réel. Elles gagnent en réactivité, en robustesse et en lisibilité, à condition que cette numérisation reste compréhensible et maintenable pour les équipes.

\subsection{Engagement DD-RSE du secteur : entre contrainte et levier stratégique}
 
La question du développement durable et de la responsabilité sociétale des entreprises (RSE) n'est pas absente du secteur du compartimentage et du désenfumage, même si elle y prend des formes spécifiques, directement liées à la nature des produits et à leurs usages.
 
\subsubsection{L'éco-conception comme enjeu de compétitivité}
 
La fabrication de solutions de compartimentage coupe-feu mobilise des matériaux dont la production et la fin de vie ont un impact environnemental mesurable : acier, aluminium, laines minérales, joints intumescents et composants de quincaillerie. Le secteur est progressivement poussé, par la réglementation comme par les donneurs d'ordre, à intégrer des démarches d'éco-conception : réduction du poids des produits, optimisation des sections structurelles, diminution des rebuts en fabrication et recours à des matériaux recyclés ou recyclables.
 
Dans ce cadre, la fiabilité de la donnée technique joue un rôle inattendu mais réel. Un configurateur bien structuré réduit les erreurs de configuration, donc les reprises, les non-conformités et les gaspillages. Il contribue indirectement à une fabrication plus sobre, en évitant que des pièces incorrectes soient produites, puis corrigées ou jetées.
 
\subsubsection{Réduction des impacts en phase de mise en œuvre et d'usage}
 
La durée de vie d'un produit de compartimentage peut s'étendre sur plusieurs décennies, car ces éléments sont intégrés au bâti. La fiabilité des performances dans le temps — maintien du classement coupe-feu, étanchéité acoustique, comportement mécanique — est donc un enjeu environnemental autant que réglementaire. Un produit bien configuré, fabriqué sans erreur et documenté avec précision a plus de chances d'être bien posé, bien entretenu et de tenir ses performances sur la durée.
 
La montée en puissance des logiques BIM dans la construction renforce cet enjeu. Les maquettes numériques intègrent désormais des données de performance, de maintenance et d'impact environnemental des produits. Les industriels qui disposent de données structurées, vérifiables et exportables vers des formats interopérables seront mieux positionnés pour répondre aux exigences croissantes des projets en matière de traçabilité et de cycle de vie.
 
\subsubsection{Engagement RSE : capital humain et transmission du savoir}
 
La RSE ne se limite pas à la dimension environnementale. Elle englobe aussi la qualité des conditions de travail, la transmission des compétences et la capacité à capitaliser le savoir de l'entreprise. Sur ce plan, les outils développés pendant cette alternance ont une dimension RSE directe : en formalisant des règles métier jusque-là portées par des experts individuels, ils réduisent la vulnérabilité de l'entreprise face aux départs ou aux transitions, et facilitent la montée en compétence de nouveaux collaborateurs.
 
Plus largement, l'entreprise qui investit dans la structuration numérique de ses données techniques crée une forme de patrimoine partageable. Ce patrimoine bénéficie aux équipes actuelles, mais il a aussi une portée sur le long terme : maintenance des produits, conformité documentaire, capacité à répondre à des audits ou à des évolutions réglementaires. Ces éléments participent à une logique de durabilité organisationnelle qui s'inscrit pleinement dans une vision RSE contemporaine.
 
\begin{table}[H]
\centering
\caption{Dimensions DD-RSE mobilisées dans le secteur et écho dans l'alternance}
\begin{tabularx}{\textwidth}{>{\bfseries}p{0.27\textwidth}X X}
\toprule
Dimension DD-RSE & Expression dans le secteur & Lien avec les missions réalisées \\
\midrule
Éco-conception & Réduction des matériaux, optimisation des gammes, recyclabilité. & Le configurateur réduit les erreurs de configuration, donc les rebuts et reprises. \\
Durabilité des performances & Maintien du classement feu, traçabilité documentaire sur la durée de vie du bâtiment. & La structuration SQL fiabilise la donnée technique à la source, réduisant les dérives de documentation. \\
Interopérabilité numérique & BIM, réutilisation des données produit sur l'ensemble du cycle de vie. & La liaison SQL--CAO prépare une continuité numérique compatible avec les démarches BIM. \\
Capital humain & Transmission des savoirs, résilience de l'organisation. & La formalisation des règles métier réduit la dépendance à des experts individuels. \\
\bottomrule
\end{tabularx}
\end{table}

\section{Positionnement de Souchier-Boullet}

\subsection{Histoire et identité de l'entreprise}

Souchier-Boullet occupe une place reconnue dans son domaine grâce à un positionnement fortement orienté vers la sécurité incendie des bâtiments, la maîtrise des ouvrants techniques et la conception de solutions répondant à des contraintes réglementaires élevées. Au-delà de la marque elle-même, l'entreprise s'inscrit dans un ensemble industriel où la spécialisation produit et la maîtrise technique constituent des marqueurs forts.

Ce qui ressort rapidement lorsqu'on découvre l'entreprise, c'est le poids de l'expérience accumulée dans les produits, les essais, les familles de solutions et la documentation associée. Une partie importante de cette richesse est explicite, via les plans, les procès-verbaux et les nomenclatures. Une autre partie reste plus diffuse, portée par les habitudes de travail, par l'expérience des ingénieurs et par les compromis progressivement stabilisés entre contraintes réglementaires, fabrication et besoins clients.

Cette double dimension, à la fois documentaire et tacite, explique en partie pourquoi la structuration numérique des données est un enjeu central pour l'entreprise. Il ne s'agit pas seulement de stocker des informations, mais d'organiser et de fiabiliser un patrimoine technique.

\subsection{Une démarche d'innovation continue}

Le service R\&D joue chez Souchier-Boullet un rôle d'interface entre plusieurs temporalités. Il doit répondre aux besoins immédiats remontés par la production ou les projets en cours, tout en préparant les évolutions plus structurantes des gammes, des modèles et des outils. Cette articulation entre court terme et moyen terme est particulièrement visible dans les missions qui m'ont été confiées.

D'un côté, les modifications de modèles 3D relevaient d'un besoin opérationnel direct : traiter des écarts, améliorer des assemblages, fiabiliser des plans. De l'autre, la construction d'un configurateur SQL portait une ambition plus transverse : capitaliser les règles métier, fluidifier le travail de plusieurs services et préparer des usages futurs, comme l'ouverture automatique du modèle 3D correspondant à une configuration.

Chez Souchier-Boullet, l'innovation ne prend donc pas uniquement la forme d'un nouveau produit. Améliorer un outil interne que tout le monde utilise au quotidien peut avoir plus d'impact qu'un développement spectaculaire mais marginal.

\section{Mon rôle au sein de l'entreprise}

\subsection{Intégration dans la dynamique R\&D}

Être intégré au service R\&D m'a placé à un carrefour : les questions de conception, de structuration de données, de validation technique et de dialogue interservices y arrivent toutes, sous une forme ou une autre. Cela oblige à comprendre le sens d'une demande avant de se précipiter sur sa résolution technique, ce qui ne m'était pas naturel au début.

L'un des apprentissages majeurs de cette expérience a été de découvrir qu'un développement utile est rarement celui qui paraît le plus élégant en théorie, mais plutôt celui qui s'insère correctement dans les pratiques de l'entreprise. Cela implique d'écouter les utilisateurs, de reformuler le besoin, de tester rapidement, puis d'itérer sans perdre la cohérence globale du système.

\subsection{Deux missions complémentaires et stratégiques}

Deux missions principales ont structuré mon parcours. La première concernait la modification et l'optimisation de modèles 3D sous Autodesk Inventor. Elle m'a permis d'entrer dans la matière produit, de comprendre la logique des assemblages et de me confronter aux contraintes de fabrication. La seconde portait sur le développement de configurateurs SQL capables d'automatiser une partie de la logique de configuration des gammes.

Ces deux missions pouvaient sembler éloignées au départ. En pratique, elles se sont révélées profondément liées. Travailler sur les modèles 3D m'a aidé à mieux comprendre la granularité des paramètres techniques. Construire le configurateur m'a ensuite obligé à formaliser cette compréhension dans une structure de données, des règles et une interface. L'intégration entre configurateur et CAO a finalement matérialisé ce lien.

\section{Problématique centrale et enjeux}

\subsection{Défis de la transformation numérique industrielle}

La transformation numérique d'une entreprise industrielle n'est pas qu'une question d'outils. Elle met en jeu des données, des habitudes, des responsabilités et des savoirs. Dans le cas étudié, plusieurs défis étaient particulièrement visibles.

Le premier défi était celui de l'hétérogénéité des sources. Les règles de configuration provenaient de procès-verbaux, d'anciens fichiers Excel, de modèles CAO et d'échanges internes. Les regrouper supposait de clarifier les notions, de trancher certains cas limites et d'accepter qu'une règle métier doive parfois être reformulée avant d'être codée.

Le deuxième défi concernait la maintenabilité. Un configurateur efficace à court terme mais impossible à faire évoluer rapidement aurait perdu une grande partie de son intérêt. Il fallait donc penser une structure suffisamment générique pour accueillir plusieurs gammes tout en restant lisible.

Le troisième défi tenait à l'interopérabilité. Relier des systèmes différents, comme une base SQL, une interface métier et un modeleur CAO, crée mécaniquement des points de fragilité. Le mémoire montre d'ailleurs que cette difficulté ne relève pas uniquement de la programmation ; elle touche aussi à la qualité des conventions de nommage, à la version des modèles et à la manière dont les responsabilités sont partagées entre équipes.

\subsection{Impact stratégique des solutions envisagées}

Les solutions développées répondent d'abord à des besoins très concrets : réduire les tâches répétitives, sécuriser la lecture des règles métier, accéder plus vite à une configuration cohérente. Mais leur portée dépasse le cadre d'un simple outil interne, car elles font circuler l'information entre des services qui, auparavant, travaillaient chacun avec leurs propres documents de référence.

À plus long terme, elles participent à la construction d'une continuité numérique plus mature. Une entreprise dont les règles métier sont explicites et exploitables peut capitaliser son expertise, étendre son configurateur à d'autres gammes et envisager de nouvelles connexions avec son écosystème numérique. Celle dont les règles restent dans la tête des experts ne le peut pas.

L'alternance présentée dans ce mémoire se situe donc à l'intersection de trois ambitions : résoudre des problèmes concrets du quotidien industriel, structurer durablement des connaissances techniques, et ouvrir des perspectives d'intégration plus avancées entre données, outils et conception.

\begin{encadretitre}{À retenir de ce chapitre}
Dans la sécurité incendie, la donnée technique est normative : une dimension ou une option renvoie toujours à un domaine de validité certifié. Souchier-Boullet dispose d'un patrimoine technique riche mais en partie dispersé entre documents, fichiers et expertise individuelle. C'est cette dispersion, plus que la technique produit, qui fonde la problématique du mémoire.
\end{encadretitre}
